% Forced finish of Q3: nonlinear Einstein-Cartan from entanglement
% at the FHHPRV (second-order) standard.
% Metric Einstein through second order is cited, not re-derived.
% The axial / Cartan sector is taken to a stop and admitted.
\documentclass[11pt]{article}
\usepackage[margin=1.05in]{geometry}
\usepackage{amsmath,amssymb,amsthm}
\usepackage{enumitem,booktabs}
\usepackage[colorlinks=true,linkcolor=blue,citecolor=blue,urlcolor=blue]{hyperref}

\newtheorem{theorem}{Theorem}[section]
\newtheorem{proposition}[theorem]{Proposition}
\newtheorem{lemma}[theorem]{Lemma}
\theoremstyle{definition}
\newtheorem{definition}[theorem]{Definition}
\newtheorem*{admission}{Admission}

\newcommand{\pP}{\textbf{[P]}}
\newcommand{\pC}{\textbf{[C]}}
\newcommand{\pO}{\textbf{[O]}}
\newcommand{\pA}{\textbf{[ADMISSION]}}
\newcommand{\kappaG}{\kappa}

\title{\textbf{Second-Order Einstein--Cartan from Entanglement:\\
The Palatini Symplectic Form and an Axial Obstruction}}
\author{Robert W.\ McGwier\\
\small CTO, Cohere Technology Group\\
\small GitHub: \texttt{n4hy}}
\date{15 August 2026 --- Version 1}

\begin{document}
\maketitle

\begin{abstract}
\noindent
The remaining open of the Emergent Gravity Program that can be forced
to a stop at a published standard is nonlinear Einstein--Cartan.
Faulkner, Haehl, Hijano, Parrikar, Rabideau and Van~Raamsdonk
\cite{FHHPRV2017} already have Einstein's equations through second
order from ball relative entropy. This note asks for the same
statement with the Palatini connection independent and torsion
algebraic in spin.

The Palatini symplectic potential is
$\theta=(1/2\kappa)\,\varepsilon_{abcd}\,e^a\wedge e^b\wedge\delta\omega^{cd}$.
It contains $\delta\omega$ and not $d(\delta e)$. Dirac coupling to
$\omega$ is algebraic and adds no $\theta$. On-shell,
$\omega=\mathring{\omega}(e)+K[s]$, the Hehl--Datta term is a
potential, and the reduced pre-symplectic form is that of
Einstein--Hilbert plus Dirac. Second-order metric dynamics therefore
collapse to FHHPRV plus a contact source $T_{\mathrm{HD}}$,
\emph{if} the only CFT data are stress-tensor one- and two-point
functions.

The axial channel does not match. CFT Fisher information for a
conserved current is nonlocal (a modular-smeared $\langle J_5 J_5\rangle$,
the same kernel that produced $\IB$). Algebraic torsion has no bulk
kinetic term and cannot reproduce that two-point function. This is
the same mismatch Paper~VI absorbed at linear order by a dictionary
improvement, now visible in the symplectic form.

\textbf{Verdict.} At the FHHPRV standard: second-order \emph{Einstein}
gravity from entanglement stands (cited). Second-order
\emph{Einstein--Cartan}, including a generated Hehl--Datta contact
as the dual of CFT axial Fisher data, is \textbf{false} without an
extra propagating bulk field, which is not Einstein--Cartan.
The Final Status nonlinear line remains \pO{} as a positive theorem
and is closed here as an obstruction in the axial sector.

No \pO{} is moved to \pP{}.
\end{abstract}

\section*{Standard of proof}

``Finish'' means a statement at the grade of \cite{FHHPRV2017}:
either the second-order EC equations (metric Einstein plus algebraic
Cartan plus Hehl--Datta) are equivalent to ball relative entropy
through $O(\varepsilon^2)$, or a specific matching identity fails.
Jacobson's entanglement equilibrium \cite{Jacobson2016} is not that
grade and is not used as a substitute.

\section{The FHHPRV skeleton}

For CFT states prepared by a Euclidean source of strength
$\varepsilon$, write
\begin{equation}
\Delta S_A-\Delta\langle H_A\rangle
=
-S(\rho_A\|\rho_A^{(0)}).
\label{eq:rel}
\end{equation}
On a family of asymptotically AdS metrics, the Hollands--Wald
identity \cite{HollandsWald2013} is
\begin{equation}
\frac{d}{d\varepsilon}\bigl(E_A^{\mathrm{grav}}-S_A^{\mathrm{grav}}\bigr)
=
\int_{\Sigma_A}\boldsymbol{\omega}
+
\int_{\Sigma_A}\mathcal{G},
\label{eq:HW}
\end{equation}
with $\boldsymbol{\omega}$ the pre-symplectic current of the gravitational
Lagrangian and $\mathcal{G}$ vanishing on-shell. If $S_A$ is
represented by HRRT and $E_A$ by the asymptotic charge, the two
left-hand sides agree and
\begin{equation}
\frac{d}{d\varepsilon}S(\rho_A\|\rho_A^{(0)})
=
\int_{\Sigma_A}\boldsymbol{\omega}+\int_{\Sigma_A}\mathcal{G}.
\label{eq:master}
\end{equation}
At first order both sides vanish except $\mathcal{G}^{(1)}$, which
forces linearized Einstein \cite{FGHMV2014}. At second order, for
sources of stress-tensor or scalar type and $\widetilde{C}_T=a_\ast$,
\cite{FHHPRV2017} compute the left-hand side from CFT two-point
functions and match it to $\boldsymbol{\omega}_{\mathrm{EH}}$, forcing
\begin{equation}
E_{ab}^{(2)}-\tfrac12 T_{ab}^{(2)}=0.
\label{eq:ein2}
\end{equation}

The only thing this note changes is the Lagrangian that defines
$\boldsymbol{\omega}$ and $\mathcal{G}$.

\section{Palatini data}

\begin{definition}[Palatini--NSM gravitational sector]
\begin{equation}
S_{\mathrm{P}}
=
\frac{1}{4\kappa}
\int\varepsilon_{abcd}\,e^a\wedge e^b\wedge R^{cd}(\omega),
\qquad
R^{ab}=d\omega^{ab}+\omega^a{}_c\wedge\omega^{cb},
\quad
\kappa=8\pi G.
\label{eq:pal}
\end{equation}
The connection is independent. Matter couples as in the New Standard
Model: fermions through $\omega$, gauge fields not through $\omega$.
\end{definition}

Varying $R^{ab}$ gives $D_\omega\delta\omega^{ab}$. After integration
by parts the boundary term is
\begin{equation}
\theta(\delta)
=
\frac{1}{2\kappa}\,\varepsilon_{abcd}\,e^a\wedge e^b\wedge\delta\omega^{cd}.
\label{eq:theta}
\end{equation}
There is no $d(\delta e)$ in $\theta$. The Dirac term
\begin{equation}
\mathcal{L}_D[\omega]
=
\frac{i}{8}\,e\,\omega_\mu^{ab}\,\bar\psi\{\gamma^\mu,\gamma_{ab}\}\psi
\label{eq:diracw}
\end{equation}
is algebraic in $\omega$ and contributes nothing to $\theta$.

The pre-symplectic $3$-form on a slice is
\begin{equation}
\boldsymbol{\omega}_{\mathrm{P}}(\delta_1,\delta_2)
=
\frac{1}{\kappa}\,\varepsilon_{abcd}\,
\delta_1 e^a\wedge e^b\wedge\delta_2\omega^{cd}
-
(1\leftrightarrow 2).
\label{eq:omegaP}
\end{equation}

\begin{proposition}[on-shell reduction]
\label{prop:red}
Write $\omega=\mathring{\omega}(e)+K$ with $K$ the contorsion.
The Cartan equation is algebraic and fixes $K$ to the Dirac spin
tensor, $K=\kappa s/2$ in the totally antisymmetric sector used by
M9.1. Substituting back produces Einstein--Hilbert plus the
Hehl--Datta potential
\begin{equation}
\mathcal{L}_{\mathrm{HD}}
=
-\frac{3\kappa}{16}\,J_5^\mu J_{5\mu}
\label{eq:HD}
\end{equation}
and no additional kinetic term for $K$ or $J_5$. The reduced
pre-symplectic form on $(e,\psi)$ is
\begin{equation}
\boldsymbol{\omega}_{\mathrm{red}}
=
\boldsymbol{\omega}_{\mathrm{EH}}(\delta e,\delta e)
+
\boldsymbol{\omega}_{\mathrm{D}}(\delta\psi,\delta\psi).
\label{eq:omred}
\end{equation}
$\mathcal{L}_{\mathrm{HD}}$ shifts the Hamiltonian. It does not
enlarge the phase space.
\end{proposition}

The metric sector of \eqref{eq:omred} is exactly the $\boldsymbol{\omega}$ of
\cite{FHHPRV2017}. Therefore:

\begin{theorem}[metric second order, cited plus reduction]
\label{thm:metric}
If the only sourced CFT operator is the stress tensor (or a scalar
primary) and $\widetilde{C}_T=a_\ast$, then \eqref{eq:master} with
$\boldsymbol{\omega}_{\mathrm{red}}$ reproduces \eqref{eq:ein2} with
$T^{(2)}$ including the Hehl--Datta stress of whatever axial
background is already present. This is FHHPRV plus
Proposition~\ref{prop:red}. It is not a new derivation of FHHPRV.
\end{theorem}

\section{The axial channel}

Now source the CFT axial current, $\delta S_{\mathrm{CFT}}=\int\beta_\mu J_5^\mu$,
as in Papers~II and~IV. At second order in $\beta$ the relative entropy
is the quantum Fisher information of $J_5$:
\begin{equation}
\delta^{(2)}S(\rho_A\|\rho_A^{(0)})
=
\int_A\!dx\int_A\!dy\,
\beta_\mu(x)\beta_\nu(y)\,
\mathcal{F}_A^{\mu\nu}(x,y),
\label{eq:fisher}
\end{equation}
where $\mathcal{F}_A$ is the modular-smeared connected two-point
function of $J_5$ --- the same object whose finite-ball integral is
$\IB$ \cite{McGwierVI}. In any unitary CFT in $d\ge 2$,
\begin{equation}
\langle J_5^\mu(x)J_5^\nu(y)\rangle
\sim
\frac{C_J}{|x-y|^{2d-2}}
\label{eq:JJ}
\end{equation}
is nonlocal.

On the gravity side, after Proposition~\ref{prop:red}, the only
axial contribution to the right-hand side of \eqref{eq:master} is
the variation of the Hehl--Datta Hamiltonian. That is ultralocal:
\begin{equation}
\delta^{(2)}H_{\mathrm{HD}}
=
-\frac{3\kappa}{8}\int J_5\cdot\delta J_5
\quad\text{with}\quad
J_5\sim\beta
\;\text{algebraically}.
\label{eq:HDloc}
\end{equation}
There is no bulk Maxwell (or other propagating) symplectic current
for an independent axial gauge field, because Einstein--Cartan does
not have one.

\begin{theorem}[axial obstruction]
\label{thm:obs}
Equations \eqref{eq:fisher}--\eqref{eq:HDloc} cannot match for a
generic smooth compactly supported $\beta$. The left-hand side of
\eqref{eq:master} is a nonlocal quadratic form (the modular Fisher
kernel). The EC right-hand side is a contact. Equality for all
$\beta$ would require
$\langle J_5(x)J_5(y)\rangle\propto\delta(x-y)$ in the ball
algebra, which contradicts \eqref{eq:JJ}.
\end{theorem}

This is the linear $\IB$ mismatch, rewritten at second order.
Paper~VI absorbed $\IB\neq 0$ into a finite dictionary improvement
at $O(\beta)$. That improvement is a redefinition of the
\emph{local} modular term. It does not install a propagating bulk
field, and it does not make \eqref{eq:fisher} ultralocal. Promoting
it to a second-order equivalence is a change of theory.

\begin{admission}
Second-order Einstein--Cartan from entanglement, in the sense
``ball relative entropy through $O(\varepsilon^2)$ $\Leftrightarrow$
nonlinear EC including generated Hehl--Datta as the dual of CFT
axial data,'' is \textbf{not a theorem}. The symplectic reduction
is clean. The matching fails in the axial channel for a structural
reason already visible at linear order. I did not produce a
Hollands--Wald charge computation that evades Theorem~\ref{thm:obs}.
I did not derive all-orders EC. I did not repair the obstruction by
adding a bulk axial gauge field: that theory is Einstein--Maxwell
(or Einstein--Proca), not Einstein--Cartan.
\end{admission}

\section{What would have counted as success}

A derivation in which
\begin{enumerate}[leftmargin=1.6em]
\item $\boldsymbol{\omega}$ is that of Palatini, not of Einstein--Maxwell;
\item $\mathcal{G}=0$ is the second-order Einstein equation \emph{and}
the algebraic Cartan equation;
\item $\mathcal{L}_{\mathrm{HD}}$ is generated, not inserted;
\item \eqref{eq:fisher} equals $\int\boldsymbol{\omega}_{\mathrm{red}}$ for
axial $\beta$,
\end{enumerate}
for all balls. Item 4 fails. Items 1--3 hold only after reducing
away from the CFT current two-point function.

\section{Verdict}

\begin{center}
\begin{tabular}{@{}p{0.40\textwidth}p{0.50\textwidth}@{}}
\toprule
Claim & Status \\
\midrule
Palatini $\theta$ is \eqref{eq:theta} & derived \\
Algebraic Dirac--$\omega$ coupling adds no $\theta$ & derived \\
On-shell reduction to EH $+$ HD potential & derived (M9.1 $+$ Prop.~\ref{prop:red}) \\
Second-order Einstein from entanglement & cited \cite{FHHPRV2017}; not re-proved \\
Second-order EC including axial Fisher $\Leftrightarrow$ HD &
\textbf{obstructed} (Thm.~\ref{thm:obs}) \\
All-orders EC on general backgrounds & still \pO{} \\
Final Status ``nonlinear regime'' as a positive & unchanged, \pO{} \\
\bottomrule
\end{tabular}
\end{center}

Forced finish: the question is answered \emph{no} at the FHHPRV
standard for the Cartan/axial sector. The metric Einstein half is
someone else's theorem. That is the stop.

\begin{thebibliography}{99}
\bibitem{FGHMV2014}
T.~Faulkner, M.~Guica, T.~Hartman, R.C.~Myers and M.~Van Raamsdonk,
JHEP 03 (2014) 051, arXiv:1312.7856.
\bibitem{FHHPRV2017}
T.~Faulkner, F.M.~Haehl, E.~Hijano, O.~Parrikar, C.~Rabideau and M.~Van Raamsdonk,
JHEP 08 (2017) 057, arXiv:1705.03026.
\bibitem{HollandsWald2013}
S.~Hollands and R.M.~Wald,
Commun.\ Math.\ Phys.\ 321 (2013) 629, arXiv:1201.0463.
\bibitem{LV2016}
N.~Lashkari and M.~Van Raamsdonk,
JHEP 04 (2016) 153, arXiv:1508.00897.
\bibitem{Jacobson2016}
T.~Jacobson,
Phys.\ Rev.\ Lett.\ 116 (2016) 201101, arXiv:1505.04753.
\bibitem{HehlDatta1971}
F.W.~Hehl and B.K.~Datta,
J.\ Math.\ Phys.\ 12 (1971) 1334.
\bibitem{Hehl1976}
F.W.~Hehl, P.~von der Heyde, G.D.~Kerlick and J.M.~Nester,
Rev.\ Mod.\ Phys.\ 48 (1976) 393.
\bibitem{McGwierVI}
R.W.~McGwier,
``Evaluation of the Universal Kinematic Integral for Condition NL,''
\url{https://github.com/n4hy/New_Model_Emergent_Gravity}.
\bibitem{McGwierFS}
R.W.~McGwier,
``Final Status of Claims,'' ibid.
\end{thebibliography}
\end{document}
